\section{Related work}
\label{sec:related-work}

Proof-carrying code places explicit evidence between a producer and a
consumer's safety policy~\citep{baseline01}. This motivates separating an
artifact from the evidence needed to accept it. Our finite-check repair
experiment and signed provenance records are not proof-carrying code in
that formal sense. Abstract interpretation and translation validation
similarly motivate explicit approximation and conversion
relations~\citep{baseline02,baseline04}, not automatic admission of a backend
or source-language proof.

Mokhov et al.\ separate scheduling from the decision whether a task needs
rebuilding~\citep{baseline07}. Our reuse obligation additionally names the
statement, environment, lifecycle, and trust conditions for transferring a
test observation. The cold comparison supplies no measured incremental-reuse
advantage over build systems.

SWE-agent studies an interface for repository navigation, editing, and program
execution~\citep{baseline08}; SWT-Bench studies test-based validation of
real-world fixes~\citep{baseline11}. They motivate treating both information
access and executable checks as parts of a repair method. Differences in tasks,
interaction, models, and acceptance criteria preclude a head-to-head performance
claim or an assertion that we reran their benchmarks. The contribution here is
a bounded composition of evidence and authority mechanisms, not a first-of-kind
claim or a new foundational algorithm.

